\documentclass[10pt]{article}
\usepackage[margin=0.55in]{geometry}
\usepackage{amsmath,amssymb}
\usepackage{hyperref}
\setlength{\parindent}{0pt}
\setlength{\parskip}{0.15em}
\setlength{\abovedisplayskip}{4pt}
\setlength{\belowdisplayskip}{4pt}
\setlength{\abovedisplayshortskip}{3pt}
\setlength{\belowdisplayshortskip}{3pt}

\begin{document}
\small
\begin{center}
{\large\bfseries A Three-Person Attendance Game: Payoff Definition}
\end{center}

\textbf{Players, states, and actions.}
There are three players, $N=\{1,2,3\}$, and every pair knows one another:
the potential contact graph is $K_3$. At the start of a day, let
$x_i\in\{0,1\}$ denote player $i$'s disease state, where $0$ means
susceptible (S) and $1$ means infected (I). The example below fixes
$x=(1,0,0)$, and all three players know this state. Each player
simultaneously chooses $a_i\in\{A,H\}$, where $A$ means attend and $H$
means stay home. A pair meets on site exactly when both attend:
\[
C_{ij}(a)=\mathbf{1}\{a_i=A\}\mathbf{1}\{a_j=A\},
\qquad i\ne j.
\]
Staying home removes all contacts in this small network. There are no
outside infections. Transmissions and recoveries use start-of-day
states; new states take effect at the end of the day.

\textbf{Probability of infection.}
For a susceptible player $i$, the number of infected neighbors who
attend alongside $i$ is
\[
m_i(a,x)=\sum_{j\in N\setminus\{i\}} x_j C_{ij}(a).
\]
Each active infected--susceptible edge independently transmits with
probability $\beta$. The probability of \emph{avoiding} infection on
every such edge is therefore
\[
\Pr(X_i'=0\mid x_i=0,x,a)
=\prod_{\substack{j\ne i\\x_jC_{ij}(a)=1}}(1-\beta)
=(1-\beta)^{m_i(a,x)}.
\]
Taking the complement gives
\[
\Pr(X_i'=1\mid x_i=0,x,a)
=1-(1-\beta)^{m_i(a,x)}.
\]
An initially infected player transmits during the day and recovers at
its end with probability $\gamma$. Thus, defining
$q_i(a,x)=\Pr(X_i'=1\mid x,a)$,
\[
q_i(a,x)=
\begin{cases}
1-(1-\beta)^{m_i(a,x)}, & x_i=0,\\
1-\gamma, & x_i=1.
\end{cases}
\]

\textbf{Realized and expected payoffs.}
Let $c_i>0$ be player $i$'s one-day cost of staying home (lost work or
social activity), and let $L_i>0$ be the cost of \emph{ending the day
infected}. Attendance has zero activity cost by normalization. If the
end-of-day infection indicator is $X_i'\in\{0,1\}$, realized payoff is
\[
u_i(a,X_i')=-c_i\mathbf{1}\{a_i=H\}-L_iX_i'.
\]
Before the random disease outcome is observed, the expected payoff is
\begin{align*}
U_i(a\mid x)
&=\mathbb{E}[u_i(a,X_i')\mid x,a]\\
&=-c_i\mathbf{1}\{a_i=H\}-L_i\mathbb{E}[X_i'\mid x,a]\\
&=-c_i\mathbf{1}\{a_i=H\}-L_i\Pr(X_i'=1\mid x,a)\\
&=\boxed{-c_i\mathbf{1}\{a_i=H\}-L_iq_i(a,x)}.
\end{align*}
The third equality holds because the expectation of a $0$--$1$
indicator equals the probability that it is $1$.

\textbf{Numerical example and best responses.}
Use $c_i=0.2$, $L_i=10$, $\beta=0.04$, and $\gamma=0.20$ for every
player; these values illustrate the game and are not empirical
estimates of the costs. With $x=(1,0,0)$, player 1's end-of-day
infection probability is $1-\gamma=0.8$, regardless of attendance:
\[
U_1(A\mid x)=-10(0.8)=-8,\qquad
U_1(H\mid x)=-0.2-10(0.8)=-8.2.
\]
Player 1 therefore prefers $A$. If player 1 chooses $A$, then each
susceptible player $i\in\{2,3\}$ faces one infected attending neighbor
when also choosing $A$. Hence
\[
\begin{aligned}
m_i(A,a_1=A,x)&=1,&
U_i(A,a_1=A\mid x)&=-10[1-(1-0.04)^1]=-0.4,\\
m_i(H,a_1=A,x)&=0,&
U_i(H,a_1=A\mid x)&=-0.2-10[1-(1-0.04)^0]=-0.2.
\end{aligned}
\]
Both susceptible players prefer $H$ when player 1 attends. Therefore
$(A,H,H)$ is a Nash equilibrium. Player 1's attendance raises others'
infection risk without raising player 1's private infection risk: this
is the externality the later policy model would need to address.

\textbf{Scope and provenance.}
This is our proposed one-day payoff, not an equation from the
telecommuting paper. The general tradeoff between distancing cost
and infection risk appears in
\href{https://doi.org/10.1371/journal.pcbi.1000793}{Reluga (2010)}.
The public infection state, perfect isolation at home, and cost of
\emph{being infected at day's end} are modeling choices to review
before game code is written.
\end{document}
